\documentclass[11pt]{article}

\newcommand{\docname}{Privacy Policy}
\newcommand{\docversion}{1.0}
\newcommand{\docdate}{\ph{EFFECTIVE DATE}}
\newcommand{\appversion}{1.0}
\newcommand{\publisher}{\ph{DEVELOPER LEGAL NAME}}

\usepackage{tactura-legal}

\hypersetup{pdftitle={Tactura — Privacy Policy},
            pdfauthor={Tactura},
            pdfsubject={Privacy Policy for the Tactura iPadOS application}}

\begin{document}
\legaltitle

This Privacy Policy explains how \publisher\ (“we”, “us”) handles information in
connection with \term{Tactura}, an iPadOS application that converts a two-dimensional
image into a 2.5D relief model suitable for 3D printing.

It is short, and it is short for a reason: \textbf{Tactura has no servers, no accounts and no
network code.} Everything the app does with your images happens on your iPad.

\begin{keybox}
{\headingfont\bfseries\fontsize{12}{14}\selectfont The one-sentence version}\par
\vspace{1.5mm}
Tactura does not collect, transmit, sell or share any personal data. Your artwork
is analysed by an artificial-intelligence model that runs entirely on your own
device, and it never leaves that device unless you deliberately export a file.
\end{keybox}

% =================================================================
\section{Summary at a glance}

\begin{plainbox}
\uifont\small
\begin{tabular}{@{}p{0.46\textwidth}p{0.47\textwidth}@{}}
\toprule
\textbf{Question} & \textbf{Answer} \\
\midrule
Do you collect personal data?            & No. \\
Do you require an account or sign-in?    & No. \\
Does the app connect to the internet?    & No. It contains no networking code. \\
Are my images uploaded anywhere?         & No. They are processed on your iPad. \\
Where does the AI run?                   & On your device, on Apple silicon. \\
Is my artwork used to train AI models?   & No. Never. \\
Do you use analytics, ads or trackers?   & No. There are no third-party SDKs. \\
Do you sell or share my data?            & No. There is nothing to sell or share. \\
Can I use the app offline?               & Yes, entirely. \\
Apple “App Privacy” declaration        & \textbf{Data Not Collected}. \\
\bottomrule
\end{tabular}
\end{plainbox}

% =================================================================
\section{The AI model runs on your device}

This is the most important thing to understand about Tactura, so it comes first.

Converting a flat image into a relief requires estimating how far every pixel of that
image is from the viewer — a task called \term{monocular depth estimation}. Most
applications that do this send your picture to a remote server, run a model there, and
send a result back. Tactura does not. The model is embedded in the application itself
and executes locally.

\subsection{What is actually running}

\begin{itemize}
  \item Tactura bundles \term{Depth Anything V2 Small}, a compact
        vision-transformer depth model published by Apple in \term{Core ML}
        format at half precision.
  \item The model file ships \emph{inside} the app you download from the App Store.
        There is no separate download step, no model server and no licence check.
  \item Inference is dispatched by Core ML to the \term{Apple Neural Engine} and the
        other processors on your iPad. Your image is read from the app's own storage,
        passed to the model in memory, and the resulting depth map is written back to
        the app's own storage.
  \item If the bundled model cannot be loaded for any reason, the app falls back to a
        purely classical image-processing method. That fallback is also entirely local
        — there is no remote path in either case.
  \item Every other stage of the conversion pipeline — image preparation, depth
        correction, region segmentation, volume construction, layer blending, mesh
        generation and file export — is likewise computed on your iPad using code
        compiled into the app.
\end{itemize}

\subsection{What this means for you in practice}

\begin{itemize}
  \item \textbf{Your images are never transmitted.} Not to us, not to a cloud
        inference provider, not to Apple, not to anyone. The application contains no
        code capable of making a network request with your content.
  \item \textbf{Nothing you make is used for training.} We do not receive your
        artwork, your prompts, your settings or your results, so none of it can be, or
        ever will be, used to train or fine-tune any machine-learning model.
  \item \textbf{The app works with no connection at all.} You can put the iPad in
        Airplane Mode and every feature — capture, crop, convert, refine, export —
        continues to work identically. This is a straightforward way to verify the
        claims on this page for yourself.
  \item \textbf{No usage telemetry is produced.} We do not know how often you open the
        app, what you convert, how long a conversion took, or whether it succeeded.
  \item \textbf{The trade-off is honest.} On-device inference uses your iPad's
        processor, battery and storage. A conversion is computationally heavy; the
        device may warm up and battery use during a conversion is noticeable. We
        consider that an acceptable price for your work staying yours.
\end{itemize}

% =================================================================
\section{Information the app can access on your device}

Tactura asks your iPad for access to three things, each only at the moment you use the
corresponding feature, and each governed by iOS's own permission system. You can
review or revoke every one of them at any time in \textbf{Settings $\rightarrow$
Privacy \& Security}, or in \textbf{Settings $\rightarrow$ Tactura}.

\subsection{Camera}
If you choose to photograph your artwork inside the app, iPadOS asks for camera
permission the first time. The camera preview is rendered locally and the still image
you capture is saved into the app's private storage as the source image for that
project. Tactura does not record video, does not capture audio, does not use the
camera in the background, and does not transmit any frame anywhere. Declining camera
access disables only the in-app capture screen; importing from your photo library
still works.

\subsection{Photo library}
If you import an existing image, Tactura presents the standard iPadOS photo picker.
That picker is provided and run by the operating system, not by us. Tactura receives
\emph{only} the single image you select. The app is not granted the ability to browse,
search, enumerate or read the rest of your photo library, and it does not request
full-library access.

\subsection{Files, when you export}
When you export a model, Tactura shows the system document picker so you can choose a
destination folder — on the iPad, in iCloud Drive, or in another provider you have
installed. iPadOS grants the app temporary, scoped permission to write into the folder
you picked, and nothing else. Tactura writes the single file you asked for and nothing
more. If you choose a destination synchronised by a cloud service, that service's own
privacy policy governs the copy you placed there; that is your choice and outside our
control.

% =================================================================
\section{What Tactura stores, and where}

All application data lives inside Tactura's sandbox on your iPad — a private area
that other apps cannot read, protected by iOS Data Protection and the device's
hardware encryption whenever the device is locked.

\begin{plainbox}
\uifont\small
\begin{tabular}{@{}p{0.26\textwidth}p{0.67\textwidth}@{}}
\toprule
\textbf{What} & \textbf{Why it exists} \\
\midrule
Project record   & The project's name, its creation and modification dates, and the
                   conversion settings you chose. \\
Source image     & The photograph or imported picture the project was made from, kept
                   so you can revisit and re-convert a project later. \\
Checkpoint file  & The intermediate result of the expensive analysis stages, cached so
                   that adjusting a slider does not re-run the AI model. It is
                   regenerable, and is deliberately excluded from device backups. \\
\bottomrule
\end{tabular}
\end{plainbox}

There is no hidden log, no crash-reporting bundle, no advertising identifier and no
usage database. We hold no copy of any of the above, because none of it is ever sent
to us.

\subsection{Retention and deletion}
Your data is retained on your device for exactly as long as you want it. Deleting a
project inside the app removes that project's folder and everything in it. Deleting
Tactura from your iPad removes all of its data. Because we store nothing, there is no
server-side copy to request the deletion of and no retention period on our side.

\subsection{Backups}
If you have enabled iCloud Backup or you back up your iPad to a computer, the project
records and source images inside the app's storage may be included in that backup, in
the same way as data from any other app. Those backups are made, encrypted and
controlled by Apple under Apple's terms and privacy policy — we have no access to
them. The large regenerable checkpoint files are explicitly marked to be excluded from
backups.

% =================================================================
\section{No third-party data collection}

Tactura contains no analytics framework, no crash-reporting service, no advertising
network, no attribution or “growth” SDK, no social-media integration and no
third-party service that receives data about you.

The app does incorporate third-party \emph{software components} — an image-processing
library, the depth model described in Section 2, and a mesh-simplification routine.
These are compiled into the application and execute locally. None of them collect
information, and none of them communicate with any network service. They are listed,
with their licences, in Tactura's Terms and Conditions.

We do not sell personal information. We do not share personal information for
cross-context behavioural advertising. We do not disclose personal information to
anyone, for any purpose, because we do not receive any.

% =================================================================
\section{Children's privacy}

Tactura is a general-audience creative tool and is suitable for all ages. It collects
no personal information from anyone, including children, and therefore collects none
from children under 13 (or the equivalent age of digital consent in your country,
which is between 13 and 16 across the European Union). There is no sign-up, no profile,
no messaging, no user-generated-content feed and no advertising in the app. Tactura is
not directed at children in a manner that would make it subject to the U.S. Children's
Online Privacy Protection Act (COPPA) as an operator collecting children's data,
because no data is collected at all.

% =================================================================
\section{Your rights under data-protection law}

Data-protection law across the world gives you rights over personal data that an
organisation holds about you. Those laws include the EU General Data Protection
Regulation (GDPR) and the UK GDPR, Indonesia's Personal Data Protection Law
(Undang-Undang No.\ 27 of 2022, \emph{Pelindungan Data Pribadi}), the California
Consumer Privacy Act as amended by the CPRA, Brazil's LGPD, Canada's PIPEDA,
Australia's Privacy Act, Japan's APPI, South Korea's PIPA, and comparable statutes
elsewhere.

In each case those rights attach to data that we, as an organisation, process about
you. \textbf{We process none.} We are not a controller or a processor of your personal
data in connection with your use of Tactura, because no personal data reaches us. There
is accordingly no international transfer of your data, no automated decision-making
carried out by us, no profiling, and no third-country recipient.

Applied to Tactura, your rights work out as follows:

\begin{plainbox}
\uifont\small
\begin{tabular}{@{}p{0.33\textwidth}p{0.60\textwidth}@{}}
\toprule
\textbf{Right} & \textbf{How it is satisfied} \\
\midrule
Access / portability  & Your data is already in your possession. Export any project as
                        a standard \texttt{.stl} or \texttt{.png} file at any time. \\
Rectification         & Edit or re-convert a project directly in the app. \\
Erasure               & Delete a project, or delete the app. Nothing survives
                        elsewhere. \\
Restriction / objection & Revoke camera or photo access in iOS Settings, or stop using
                        the app. \\
Withdraw consent      & Revoke the relevant iOS permission at any time; the app
                        continues to function without it. \\
Non-discrimination    & Every feature is available to every user. There is no paid
                        tier that trades features for data. \\
Complaint             & You may complain to your national supervisory authority — in
                        Indonesia, the authority designated under UU PDP; in the EU or
                        UK, your local data-protection authority. \\
\bottomrule
\end{tabular}
\end{plainbox}

If you would like this confirmed in writing for your own records, write to us at the
address in Section 10 and we will respond.

\subsection{A note for California residents}
In the twelve months preceding the effective date of this policy, we have not
collected, sold or shared any category of personal information listed in the CCPA. We
have no “Do Not Sell or Share My Personal Information” link because there is nothing
to opt out of.

% =================================================================
\section{Security}

Because Tactura keeps everything local, its security posture rests on your device's own
protections, which are substantial:

\begin{itemize}
  \item Application data is stored inside the iOS sandbox, which prevents other
        applications from reading it.
  \item Files are protected by iOS Data Protection and are encrypted at rest by the
        device's hardware while the device is locked.
  \item The app is distributed exclusively through the App Store and is code-signed;
        iPadOS verifies that signature before it will run.
  \item There is no attack surface in the form of an account, a password, a session
        token, an API key or a server, because none of those exist.
\end{itemize}

The practical advice that follows from this is ordinary but worth stating: protect your
iPad with a passcode or Face~ID, keep iPadOS updated, and remember that anyone with
access to your unlocked device has access to your projects.

% =================================================================
\section{Changes to this policy}

If Tactura ever gains a feature that changes the answers above — for example an
optional cloud backup, or crash reporting — we will update this policy before that
feature ships, revise the “Last updated” date at the top, and describe the change in
the App Store release notes for the version concerned. Material changes that would
involve collecting personal data will be presented to you in the app, and will require
your explicit opt-in. We will not quietly begin collecting data under an old policy.

Previous versions of this policy remain available on request.

% =================================================================
\section{Contact}

Questions about this policy, or about privacy in Tactura generally, are welcome.

\begin{plainbox}
\uifont
\textbf{Data controller / publisher}\quad \publisher \\[3pt]
\textbf{Email}\quad \ph{CONTACT EMAIL} \\[3pt]
\textbf{Postal address}\quad \ph{POSTAL ADDRESS} \\[3pt]
\textbf{EU / UK representative}\quad \ph{IF APPLICABLE — OTHERWISE DELETE THIS LINE} \\[3pt]
\textbf{Support}\quad \ph{SUPPORT URL}
\end{plainbox}

We aim to reply to any privacy enquiry within 30 days, which is the response window
required by the GDPR and by most comparable statutes.

\vspace{4mm}
{\color{tactRule}\hrule height 0.6pt}
\vspace{2mm}
{\uifont\footnotesize\color{tactMuted}
This policy covers the Tactura application only. It does not cover the App Store
itself, iCloud, or any third-party service you may choose to export a file into —
each of those is governed by its own provider's privacy policy.\par}

\end{document}
